% !TeX root = ../thuthesis-example.tex

\chapter{栖息无人机的控制实验与讨论}
本章围绕微型无人机栖息任务的控制实验与系统性能验证展开，旨在基于前述硬件平台与控制框架，对所设计的视觉伺服控制系统、姿态稳定控制器以及 VSERF 柔性黏附栖息方案进行实验验证与性能分析，进一步评估系统在真实环境下的控制效果、鲁棒性与工程可行性。

首先，将对实验平台及测试流程进行介绍，包括地面站视觉处理系统、无人机机载控制架构、AprilTag 目标环境以及 VSERF 黏附执行器的实验部署方式，并说明实验过程中涉及的定义、控制频率与关键参数设置，以保证实验结果的可重复性与分析一致性。

随后，将围绕无人机姿态控制性能与视觉伺服能力开展基础实验，通过悬停稳定性测试、目标跟踪与自主接近实验，对姿态内环控制器的稳定调节能力以及视觉外环在目标定位与误差收敛过程中的动态表现进行分析。同时，结合实验数据对视觉观测误差、通信时延及受限机载感知条件下系统闭环性能的影响进行讨论，以验证第三章理论分析与仿真结果在实际系统中的适用性。

在此基础上，将进一步开展无人机自主栖息实验，对系统从目标识别、轨迹接近、姿态调整到最终附着过程中的整体任务执行能力进行验证。此外，还将针对实验过程中出现的误差来源与失效模式进行讨论，包括视觉尺度退化、姿态扰动、气流干扰以及结构柔性带来的不确定性，以总结系统在复杂实际环境中的性能边界与改进方向。

\section{实验平台硬件环境搭建}
为验证无人机自主栖息过程的可行性并实现系统闭环控制，本文搭建了如图~\ref{experiment_setup}所示的实验平台。该实验平台主要由起飞平台、栖息平台及辅助照明系统等组成，用于为无人机自主起飞、视觉感知、目标接近与附着实验提供稳定的测试环境。

\subsection{栖息场景搭建}
其中，栖息平台固定安装于三轴手动滑台上，可便捷调节其空间位置及相对高度，从而实现目标平面与地面距离的灵活调整，以满足不同实验条件下无人机接近与栖息测试需求。滑轨距离地面的高度可以从\SI{0}{cm}至\SI{80}{cm}调节。

\begin{figure}[htbp]
  \centering
  \includegraphics[width = 0.8\textwidth]{experiment_setup.png}
  \caption{栖息实验场景搭建}\label{experiment_setup}
\end{figure}

\subsection{无人机及地面站系统硬件搭建}
为实现无人机视觉信息接收与控制指令生成，系统搭建了地面站装置，其主要设备如图~\ref{ground_station_set}所示。地面站主要包括视频接收高频头、图像传输接收设备以及控制计算终端等，用于完成机载图像接收、视觉信息处理及控制指令发送等功能。

\begin{figure}[htbp]
  \centering
  \includegraphics[width = 0.8\textwidth]{gs_set.jpg}
  \caption{地面站系统硬件选型}\label{ground_station_set}
\end{figure}

实验过程中，将预先打印的 AprilTag 标签固定于光源下方，其主要目的在于利用光源作为栖息目标载体，模拟无人机在真实环境中向路灯等高位设施执行自主栖息任务的应用场景。实验预期中，无人机在视觉伺服控制作用下完成对 AprilTag 的识别与相对位姿估计，并逐步接近目标区域，最终执行附着动作并稳定栖息于光源下方。

依据章节~\ref{Drone_system_design}所示的系统设计方案，完成无人机平台的硬件搭建。系统所采用的主要硬件组成及参数信息如表~\ref{Drone_part_list}所示。

\begin{table}[htbp]
    \centering
    \caption{无人机系统主要零件清单}
    \label{Drone_part_list}
    \begin{tabular}{cccc}
        \toprule
        零件名称 & 型号 & 数量 & 备注 \\
        \midrule
        机架 & N/A & 1 & 碳纤维加工件 \\
        相机电池架 & N/A & 1 & 3D 打印件 \\
        无刷电机 & T-HOBBY M0802 & 4 & -- \\
        螺旋桨 & T-MOTOR M12199 & 4 & -- \\
        飞控主板 & F4 AIO 1S & 1 & STM32F411，集成电调 \\
        摄像头 & FOXEER 2 MICRO 1.7 & 1 & 分辨率 1200TVL \\
        电池 & DAI WANG GAU 1S & 1 & 容量 \SI{550}{mAh} \\
        \bottomrule
    \end{tabular}
\end{table}

\section{无人机控制器验证}
本节对所设计的无人机控制器进行系统性验证，主要从仿真误差分析与实机飞行实验两个层面展开。在仿真层面，通过记录无人机在飞行过程中的三轴误差（$X$、$Y$、$Z$），并将其映射至空间轨迹，以评估控制器在不同参数设置下的收敛性与稳定性。在实机层面，基于视觉伺服系统与AprilTag定位方法，对无人机在真实环境中的自主接近与栖息过程进行多次重复实验，从而验证控制器在存在外界扰动条件下的鲁棒性与实际可行性。整体验证流程如图~\ref{fig:Error_traj}、图~\ref{fig:controller_verify}及图~\ref{fig:Path_verify}所示。

\subsection{控制器数据采集}
设计程序用于记录无人机飞行过程中的状态数据，主要包括 $X$、$Y$、$Z$ 三个方向的误差量。通过简单的坐标变换，该误差可以进一步映射为无人机在空间坐标系下的相对位置，从而用于分析其运动轨迹与控制性能。如图~\ref{fig:Error_traj}所示，六组实验分别对应六种不同的控制器参数或工作状态。

\begin{figure}[htbp]
    \centering

    \begin{subfigure}[b]{0.48\textwidth}
        \centering
        \includegraphics[width=\textwidth]{Error_Trajectory_1.png}
        \caption{轻微振荡但整体收敛}
    \end{subfigure}
    \hfill
    \begin{subfigure}[b]{0.48\textwidth}
        \centering
        \includegraphics[width=\textwidth]{Error_Trajectory_2.png}
        \caption{较大幅度振荡但仍趋于收敛}
    \end{subfigure}

    \vspace{1.0em}

    \begin{subfigure}[b]{0.48\textwidth}
        \centering
        \includegraphics[width=\textwidth]{Error_Trajectory_3.png}
        \caption{收敛性能较优}
    \end{subfigure}
    \hfill
    \begin{subfigure}[b]{0.48\textwidth}
        \centering
        \includegraphics[width=\textwidth]{Error_Trajectory_4.png}
        \caption{高度方向出现明显发散}
    \end{subfigure}

    \vspace{1.0em}

    \begin{subfigure}[b]{0.48\textwidth}
        \centering
        \includegraphics[width=\textwidth]{Error_Trajectory_5.png}
        \caption{存在单向漂移，未形成闭环收敛}
    \end{subfigure}
    \hfill
    \begin{subfigure}[b]{0.48\textwidth}
        \centering
        \includegraphics[width=\textwidth]{Error_Trajectory_6.png}
        \caption{几乎无有效控制响应}
    \end{subfigure}

    \caption{无人机误差轨迹及不同控制器性能对比}
    \label{fig:Error_traj}
\end{figure}

从图中可以观察到，不同控制器参数下系统呈现出显著不同的动态特性。实验1–3整体表现为误差逐渐收敛，但伴随不同程度的振荡，其中实验3具有最优的收敛性能与最小的稳态波动。实验4在高度方向出现明显发散，表明纵向通道控制增益可能存在不匹配或耦合未充分补偿的问题。实验5表现为明显的单向漂移，说明系统未形成有效闭环约束或存在积分饱和现象。实验6则几乎未产生有效控制输出，可能对应控制器失效或执行器饱和状态。

上述结果表明，该误差轨迹可作为评估控制器稳定性与收敛性能的直观指标，能够有效区分不同控制策略在三维空间中的动态表现。

\subsection{实机测试结果}
将无人机搭建完毕，放入测试场地进行无VSERF器件测试，以验证仿真及数值计算中控制器的效果。无人机从起飞，能够自主寻找AprilTag，并且进行导航、轨迹规划，直到无人机与待黏附表面发生碰撞。在无VSERF器件的情况下，无人机发生碰撞后会弹开，无黏附效果。实验过程如图~\ref{fig:controller_verify}所示。

\begin{figure}[htbp]
    \centering

    \begin{subfigure}[b]{0.3\textwidth}
        \centering
        \includegraphics[width=\textwidth]{verify_01.png}
        \caption{起飞阶段}
    \end{subfigure}
    \hfill
    \begin{subfigure}[b]{0.3\textwidth}
        \centering
        \includegraphics[width=\textwidth]{verify_02.png}
        \caption{视觉对准}
    \end{subfigure}
    \hfill
    \begin{subfigure}[b]{0.3\textwidth}
        \centering
        \includegraphics[width=\textwidth]{verify_03.png}
        \caption{视觉对准}
    \end{subfigure}
    
    \vspace{1.0em}

    \begin{subfigure}[b]{0.3\textwidth}
        \centering
        \includegraphics[width=\textwidth]{verify_04.png}
        \caption{视觉对准}
    \end{subfigure}
    \hfill
    \begin{subfigure}[b]{0.3\textwidth}
        \centering
        \includegraphics[width=\textwidth]{verify_05.png}
        \caption{视觉对准}
    \end{subfigure}
    \hfill
    \begin{subfigure}[b]{0.3\textwidth}
        \centering
        \includegraphics[width=\textwidth]{verify_06.png}
        \caption{视觉对准}
    \end{subfigure}

    \vspace{1.0em}

    \begin{subfigure}[b]{0.3\textwidth}
        \centering
        \includegraphics[width=\textwidth]{verify_07.png}
        \caption{上升接近（无视觉引导）}
    \end{subfigure}
    \hfill
    \begin{subfigure}[b]{0.3\textwidth}
        \centering
        \includegraphics[width=\textwidth]{verify_08.png}
        \caption{上升接近（无视觉引导）}
    \end{subfigure}
    \hfill
    \begin{subfigure}[b]{0.3\textwidth}
        \centering
        \includegraphics[width=\textwidth]{verify_09.png}
        \caption{碰撞后弹开}
    \end{subfigure}

    \caption{无人机无VSERF器件的情况下控制器的验证}
    \label{fig:controller_verify}
\end{figure}

随后，对无人机自主接近与栖息过程进行了多次重复实验，以验证控制器在实际环境中的稳定性与抗扰动能力。由于实验过程中环境条件难以完全保持一致，例如微小气流扰动以及初始姿态误差等因素的影响，无人机在接近目标并完成黏附前的飞行路径存在一定差异。为更直观地展示无人机在不同实验条件下的运动过程，本文采用图像堆栈方式对三次实验的飞行过程进行可视化，实验结果如图~\ref{fig:Path_verify}所示。

\begin{figure}[htbp]
  \centering
  \begin{subfigure}[b]{0.7\textwidth}
  \centering
  \includegraphics[width=\textwidth]{path_verify_01.png}
  \caption{实验1}
  \end{subfigure}

  \vspace{1.0em}

  \begin{subfigure}[b]{0.7\textwidth}
  \centering
  \includegraphics[width=\textwidth]{path_verify_02.png}
  \caption{实验2}
  \end{subfigure}

  \vspace{1.0em}

  \begin{subfigure}[b]{0.7\textwidth}
  \centering
  \includegraphics[width=\textwidth]{path_verify_03.png}
  \caption{实验3}
  \end{subfigure}

  \caption{三次无人机自主接近实验的飞行路径图像堆栈结果}
  \label{fig:Path_verify}

\end{figure}

实验结果表明，尽管受到外部扰动影响，无人机在多次实验中均能够稳定完成目标接近过程，并最终沿期望方向飞向目标区域。虽然不同实验中的具体飞行轨迹存在一定偏差，但整体路径均表现出良好的收敛性与一致性，说明所设计的视觉伺服控制器具备一定的抗扰动能力与路径引导能力，能够有效完成黏附前的自主接近与路径跟踪任务。



\section{VSERF器件验证}
依据课题组所提出的方法制备 VSERF 柔性黏附器件~\cite{pan2022stretchable}，并将其集成至当前无人机平台，完成相关驱动电路连接后置于栖息实验环境中。为验证 VSERF 器件实际承载能力及其与无人机控制系统之间的通信链路可靠性，首先开展人工辅助栖息实验。

实验过程中，通过人为施加接触预载荷，使无人机与栖息目标表面充分接触，并将无人机悬挂于目标下方；随后启动 VSERF 黏附器件，其实验结果如图~\ref{fig:hand_stick}所示。可以观察到，在器件工作状态下，无人机能够稳定悬挂于目标表面下方，未出现明显脱附或姿态失稳现象。实验结果表明，所制备的 VSERF 柔性黏附器件具备满足无人机栖息任务需求的承载能力，同时验证了高压驱动电路与控制通信链路工作的有效性，为后续自主栖息实验提供了系统基础。

\begin{figure}[htbp]
  \centering
  \begin{subfigure}[b]{0.48\textwidth}
  \centering
  \includegraphics[width=\textwidth]{hand_stick_01.png}
  \end{subfigure}
  \hfill
  \begin{subfigure}[b]{0.48\textwidth}
  \centering
  \includegraphics[width=\textwidth]{hand_stick_02.png}
  \end{subfigure}

  \caption{人工辅助栖息实验结果及器件验证}
  \label{fig:hand_stick}
\end{figure}

\section{无人机栖息实验}

\section{本章小结}